\section{Main theorem and the four-phase case}

\begin{theorem}[Dihedral receipt parity]
Let (n\geq3), let (A) and (B) be involutions, and suppose
(H=AB) has order (2n). Put
\[
\comm=[A,B],
\qquad
\receipt=H^n.
\]
Then:
\begin{enumerate}
\item (langle A,B\rangle) is dihedral of order (4n);
\item (H^{-1}=BA), (comm=H^2), and (|\comm|=n);
\item (
eceipt) is the unique order-two element of (langle H\rangle);
\item if (n) is even, then (
eceipt=\comm^{n/2});
\item if (n) is odd, then
\(langle H\rangle=\langle\comm\rangle\times\langle\receipt\rangle\)
and (H=\comm^{(n+1)/2}\receipt);
\item in the natural encounter quotient (q:\langle A,B\rangle\to V_4),
(q(\receipt)=1) if and only if (n) is even.
\end{enumerate}
\end{theorem}

\begin{proof}
The dihedral normal-form argument gives statements 1 and 2. The cyclic
subgroup of order (2n) has the unique involution (H^n), proving statement
3. The congruence (2k\equiv n\pmod{2n}) is solvable exactly when (n) is
even, yielding statement 4. For odd (n), the subgroups generated by (H^2)
and (H^n) have coprime orders (n) and (2), trivial intersection, and
product of order (2n); the displayed exponent identity reconstructs (H),
proving statement 5. Finally (q(H)=ab), so
(q(H^n)=(ab)^n), proving statement 6.
\end{proof}

For (n=4), (H) has order eight and the completion has order sixteen.
The even law gives
\[
\comm=H^2,
\qquad
\receipt=\comm^2=H^4.
\]
This is precisely the earlier four-phase receipt tower. The new theorem shows
both why it closes and which part of that closure depends on the evenness of
four.

